\documentclass[11pt]{article}
\usepackage{amsmath,amssymb}
\usepackage[margin=1in]{geometry}
\newcommand{\diag}{\operatorname{diag}}
\newcommand{\sgn}{\operatorname{sgn}}
\title{The fixed testing span and a nonuniform clipped Gram at one C1 contact}
\author{}\date{}
\begin{document}\maketitle

\paragraph{Scope.}
Consider only the bounded-$z$, positive-angular-coefficient C1 frame of
\texttt{latex/certificates/c1\_short.tex}, equations P1 and SP3--SP6
and the twelve-root conclusion at lines 7672--7687.
The exact clipped matrix considered below is the matrix $G$ defined in
\texttt{analysis/exact-lift-2026-10-23/fixed\_box\_c1.tex},
Section~3, by $g^{\mathsf T}Gg=g^{\mathsf T}D_*g+\mathcal E$.
Its fixed-box realization is an existing result. The new conclusion here
concerns extending this particular matrix beyond that box; it is not an
obstruction to its fixed function span or to a different Gram matrix.

\paragraph{The fixed testing span before the coefficient argument.}
Use nonnegative testing angles $u,w$. The exact angle gauges are
\[
 (u,w)=(|u_s|,|w_s|),\quad
 (W,X,Y,Z)=(W_s,\sgn_+(w_s)X_s,\sgn_+(u_s)Y_s,
             \sgn_+(u_s)\sgn_+(w_s)Z_s),
\]
where $\sgn_+(0)=1$. Put
\[
 \omega=\sgn(W)\sqrt{uw\max(WZ,0)},\qquad \omega=0\text{ at }W=0.
\]
The twelve fixed pure-testing functions are
\[
 \mathcal V=\{X,Y,Z,W,uZ,wZ,uW,uwW,u^2W,u^{3/2}W,u^{1/2}W,\omega\}.
 \tag{1}
\]
Indeed, write $q=\sqrt z$, $k=\sqrt x$, $Q=1-b$ and use the
source quantities
\[
 d_G=4D/\sigma^2,\quad \rho=\Lambda x/\sigma,\quad q_*=\Lambda k,
 \qquad \Lambda>0.
\]
The roots are source-linear combinations of (1), since
\[
\begin{gathered}
 N_x=X+uhZ,\quad N_y=(Y-bwZ)/q,\quad
 R=wZ/q-uqW,\quad \Gamma_*=\sqrt{d_G}(\sigma Z-uwW),\\
 J_C=N_y+bR,\qquad J_\ell=N_y-QR,\\
 A_\theta=q(u^2-k^2)W,\qquad T_\theta=kq(u-k)W,\\
 B_1=\sqrt{\Lambda\sigma}\,uZ,\quad
 B_2=\sqrt\rho\,uwW,\quad B_3=\sqrt\Lambda\,xqW,\quad
 B_4=\sqrt{\Lambda x}\,wZ/q,\\
 d_i=B_i-\sqrt{\Lambda x}\,\omega,\qquad
 U_*=\sqrt{q_*}\,q(u^{3/2}-ku^{1/2})W,\qquad
 K_*=\sqrt{q_*k}\,q(u-k)W.
\end{gathered} \tag{2}
\]
Here the zero $D=0$ simply sets $\Gamma_*=0$; no inverse of $D$ is used.
In particular $u-k$ is outside every radical. There is no
source-dependent testing branch point in (2). The testing branch loci
$u=0$ and $WZ=0$ in (1), and the angle sign partition, are fixed.
All divisions and radicals involving only $z,k,\sigma,D,\Lambda$
are source coefficients.

For unit states, the physical slack is $F_w/[(1+u^2)(1+w^2)]$.
Thus its roots acquire only the fixed pure-testing factor
$[(1+u^2)(1+w^2)]^{-1/2}$. One may either restrict the state to the
unit sphere or, away from the zero state, add its inverse norm as a
fixed testing normalization. The homogeneous identity itself includes
the zero state.

What remains after (2) is a coefficient Gram problem. The generic
TF4 operation in \texttt{03\_tensor.tex}, lines 58--72, introduces
testing gates $\sqrt{M_j\pm y_j}$ with thresholds chosen for the
individual bounded tensor. A moving threshold is not absorbed by
a source coefficient: the functions $\sqrt{M+t}$ for distinct positive
$M$ are linearly independent on a common interval about zero, as follows
by differentiation and a Vandermonde matrix. A fixed replacement
therefore needs an additional uniform allocation or a new identity.

\paragraph{What was already evaluated.}
The earlier note \texttt{fixed\_box\_c1.tex} and
\texttt{c1\_retained\_feature\_union\_body.tex} give fixed gates on
\[
 0<x,\sigma\le2^{-160},\quad 0<z\le1,\quad
 0<b<1,\quad bz<1,\quad 0<h<1,\quad
 E,M,J>0,\quad D,\beta\ge0,\quad \Lambda>0,
 \qquad 0\le u,w\le2^{-100}.
 \tag{3}
\]
There is no upper bound on $h/\sigma$, and both event-zero faces are
included. With
\[
 \mathcal B=\{w,w^2,\sqrt u,u,uw,u^2,u^2w\},\qquad
 \gamma=(7\,2^{38})^{-1},
\]
the main library is
$\mathcal V\{1,\sqrt{\gamma+b},\sqrt{\gamma-b}:b\in\mathcal B\}$,
with 180 labelled functions. The retained squares enlarge it to 220.
The subsequent published local-order669 construction is not a new
consequence of this note. In particular (3) already answers the fixed
library question affirmatively on one numerical source box.

\paragraph{Exact native source and a family beyond that box.}
For completeness the source formulas used in the calculation are
\[
\begin{gathered}
 r=bz,\quad \theta=1-h,\quad e=xh/\sigma,\quad E=1-e,\quad K=1+x,\\
 R_s=1+xr,\quad M=R_s-xz\theta,\quad
 J=1+xh-\sigma(1-r),\\
 P=\sigma+K+\sigma z(1+bx+2xh)
       +\sigma bx^2hz^2+x^2h\theta z,\\
 p=\frac{z(1+xh)E+(1-r)(1+xhz)}{2P},\\
 l=\frac{(1+\sigma r)\{eR_s(1+\sigma z)+ME\}}{2P},\qquad
 \kappa=\frac{K(1+\sigma r)(1+\sigma z)}{2MP},\\
 m=\kappa x^2zh/\sigma,\quad
 C=\frac{Kz(1+\sigma r)E}{2(1-r)P},\quad
 A=p/r,\quad \alpha=l/(zQ),\\
 D=\frac{K(1+\sigma z)(1+\sigma r)-JP}{2JP},\qquad
 \beta=\frac{K(1+\sigma r)(1+xhz)-\theta P}{2\theta P}.
\end{gathered} \tag{4}
\]
Fix any real $R>4$ and let $\tau\downarrow0$ in
\[
 z=R,\quad b=\frac1{2R},\quad r=\frac12,\quad
 h=\frac34,\quad \theta=\frac14,\quad
 \sigma=\tau,\quad x=\tau^3.
 \tag{5}
\]
This is inside the original bounded-$z$ source type when $\tau$ is
infinitesimal, for every fixed finite $R$. No common upper bound on
all such real $R$ is part of that type.

The six displayed event weights are strictly positive for sufficiently
small positive $\tau$; the two omitted C1 weights remain zero.
The threshold may depend on $R$.
Indeed $0<b,r,h,\theta<1$, $E=1-3\tau^2/4>0$,
$M=1+(1/2-R/4)\tau^3\to1$, and
$J=1-\tau/2+3\tau^3/4\to1$. Also $P\to1$ and
\[
 H_0=1+zQ=R+\tfrac12,\quad
 p\longrightarrow H_0/2,\quad l,\kappa\longrightarrow\tfrac12,
 \quad \beta\longrightarrow\tfrac32,\quad
 \frac D{\tau^2}\longrightarrow\frac{H_0}{2}>0.
 \tag{6}
\]
The last limit also follows without expanding (4): put
$B_s=P+xz(J-\sigma)$ and use the exact source identities
\[
 2\theta P\beta=hB_s-\sigma H_0,\qquad
 2P(JD+x\theta\beta)=\sigma^2\mathcal H,
 \quad \mathcal H\longrightarrow H_0.
\]
Thus $m,C,A,\alpha,D,\beta$ are positive in (5) for small $\tau$.

The angular quantities are
\[
 \Lambda=\frac{8Dh+4\beta\theta^2+2h^2}{\sigma}
                  -4bp\sigma z,\quad
 K_A=4\kappa h/\sigma-\Lambda,\quad c=\max(0,1-K_A).
\]
Direct substitution, or first-order expansion of (4), gives
\[
 \tau\Lambda\longrightarrow\tfrac32,\quad
 K_A\longrightarrow\tfrac{R-1}{2}>1,\quad
 \rho\longrightarrow0,\quad \Lambda x\longrightarrow0.
 \tag{7}
\]
Hence $\Lambda>0$ and $c=0$ identically for sufficiently small $\tau$.

\paragraph{The two-by-two limit.}
The literal clipped matrix uses
\[
 g=(N_x,J_C,J_\ell,\Gamma_*,A_\theta,T_\theta,d_1,\ldots,d_4,U_*,K_*),
 \qquad
 D_*=\diag(2,1,1,1,\tfrac12,\tfrac12,
                    \tfrac12,\tfrac12,\tfrac12,\tfrac12,1,\tfrac12).
\]
Write $g_s=h-Q$. The entries involving $J_C,A_\theta$ are determined by
\[
\begin{split}
 A_y&=8(p-l)+4z\{(1+2m)g_s+2Dh+2p\sigma\},\\
 \chi&=8(pb+lQ)-4h+4eh(2D+1)-8Dh\\
 &\hspace{8mm}+4z\{4Dbh+4pb\sigma+(1+2m)g_s^2+2\beta\theta^2\}.
\end{split} \tag{8}
\]
The cross terms are $uA_yN_yA_\theta+\chi uA_\theta R$.
Since $N_y=QJ_C+bJ_\ell$ and $R=J_C-J_\ell$, the symmetric
$J_C,A_\theta$ entry is $u(QA_y+\chi)/2$. Equations (5)--(6) give
\[
 g_s\to-\tfrac14+\frac1{2R},\qquad
 A_y\to3R,\qquad \chi\to R+2,\qquad
 (QA_y+\chi)/2\to2R+\tfrac14. \tag{9}
\]
At $w=0$ the radial coefficients have limits
\[
 E_{yy}\to2Ru^2,\quad
 E_{RR}\to(R/2-1/2+1/(2R))u^2,\quad
 E_{yR}\to(-R+2)u^2.
\]
Consequently
\[
 Q^2E_{yy}+E_{RR}+QE_{yR}\longrightarrow\tfrac32Ru^2.
\]
The remaining radial correction $-\Lambda xR^2/2$ tends to zero.
The second-normal radial correction also tends to zero; its coefficient
is $R_0\rho\sigma z/S$, where
$R_0\to0$ and $S=d_G+\rho\to2H_0>0$.
No other term contributes to this principal block. Because $c=0$,
its $A_\theta$ diagonal is exactly $1/2$. Thus, for each fixed $u$,
\[
 G[\{J_C,A_\theta\}]\longrightarrow
 \begin{pmatrix}
  1+\frac32Ru^2 &(2R+\frac14)u\\
  (2R+\frac14)u&\frac12
 \end{pmatrix},
 \quad
 \det\longrightarrow
 \frac12-\left(4R^2+\frac R4+\frac1{16}\right)u^2.
 \tag{10}
\]
For any proposed real cutoff $\delta>0$, fix
$0<u\le\min(\delta,1/2)$ and then choose a finite $R>4$ making
(10) negative. Continuity makes the actual determinant negative for
all sufficiently small positive $\tau$. Therefore this exact clipped
matrix is not positive semidefinite on any common real testing box
$0\le u,w\le\delta$ over all sources (5).

\paragraph{Exact remaining dependence.}
The source-to-root map (2) has a fixed testing span. What fails
uniformly is this clipped matrix's allocation after positive source
excesses have been retained separately: the mixed coefficient grows
as $2Ru$ while its clipped angular pivot stays $1/2$.
For each fixed finite $R$, independent infinitesimal $u,w$ still make
that entry infinitesimal, so (10) does not contradict the qualitative
strict-residue argument. It shows why that argument does not supply
one numerical cutoff as $R$ varies over all finite values.

A different source-dependent Gram in the same testing span, a
completion retaining the omitted positive excesses, or a further
source normalization might remove this dependence. Neither (10)
nor the independence of moving TF4 gates excludes those possibilities.
In particular this note proves no failure of a finite library for the
whole native C1 family, no exclusion from the forty-function
ramified-completion library, and no global coverage statement.
\end{document}

